\documentclass[a4paper,11pt]{article}
\usepackage[margin=22mm,headheight=16pt,footskip=10mm]{geometry}
\usepackage{fontspec,xeCJK}
\setmainfont{Times New Roman}
\setsansfont{Arial}
\setmonofont{Consolas}[Scale=0.88]
\setCJKmainfont{PMingLiU}[BoldFont={Microsoft JhengHei Bold},ItalicFont={PMingLiU},BoldItalicFont={Microsoft JhengHei Bold}]
\setCJKsansfont{Microsoft JhengHei}[BoldFont={Microsoft JhengHei Bold}]
\setCJKmonofont{PMingLiU}
\usepackage{amsmath,amssymb,amsthm}
\usepackage{booktabs,tabularx,longtable,array}
\usepackage{enumitem,setspace,titlesec,fancyhdr}
\usepackage{needspace}
\usepackage[hidelinks,unicode,bookmarksnumbered]{hyperref}
\usepackage{xurl}
\setstretch{1.22}
\setlength{\parindent}{2em}
\setlength{\parskip}{3pt}
\setlength{\emergencystretch}{2em}
\clubpenalty=10000
\widowpenalty=10000
\renewcommand{\tablename}{表}
\renewcommand{\refname}{參考文獻}
\setlength{\tabcolsep}{5pt}
\renewcommand{\arraystretch}{1.15}
\setlist[itemize]{topsep=3pt,itemsep=2pt,leftmargin=2em}
\setlist[enumerate]{topsep=3pt,itemsep=2pt,leftmargin=2em}
\titleformat{\section}{\Large\sffamily\bfseries}{\thesection}{.7em}{}
\titleformat{\subsection}{\large\sffamily\bfseries}{\thesubsection}{.7em}{}
\titlespacing*{\section}{0pt}{14pt}{7pt}
\titlespacing*{\subsection}{0pt}{10pt}{5pt}
\pagestyle{fancy}\fancyhf{}
\fancyhead[L]{\small 大稻埕人流預測與列車調度}
\fancyhead[R]{\small 理論、假設與證明}
\fancyfoot[C]{\thepage}
\renewcommand{\headrulewidth}{0.3pt}
\newtheoremstyle{chinese}{6pt}{6pt}{\normalfont}{}{\bfseries}{。}{0.5em}{}
\theoremstyle{chinese}
\newtheorem{theorem}{定理}[section]
\newtheorem{proposition}[theorem]{命題}
\newtheorem{lemma}[theorem]{引理}
\newtheorem{definition}[theorem]{定義}
\newtheorem{remark}[theorem]{說明}
\renewcommand{\proofname}{證明}
\numberwithin{equation}{section}
\allowdisplaybreaks[1]
\newcommand{\R}{\mathbb R}
\newcommand{\one}{\mathbf 1}
\newcommand{\E}{\mathbb E}
\newcommand{\Var}{\operatorname{Var}}
\newcommand{\argmin}{\operatorname*{arg\,min}}
\newcommand{\diag}{\operatorname{diag}}
\newcommand{\rank}{\operatorname{rank}}
\newcommand{\pos}[1]{\left[#1\right]_+}
\newcommand{\norm}[1]{\left\lVert#1\right\rVert}
\hypersetup{pdftitle={大稻埕人流預測與列車調度：理論、假設及最佳解存在性},
pdfauthor={專案方法文件},pdfsubject={Ridge、GAM、XGBoost、融合及批次候車最佳化}}
\begin{document}
\begin{center}
{\LARGE\sffamily\bfseries 大稻埕人流預測與列車調度}\\[5pt]
{\Large\sffamily\bfseries 理論、假設及最佳解存在性}\\[8pt]
專案方法與數學證明文件\\
2026 年 10 月 18 日
\end{center}
\vspace{4pt}
\noindent\textbf{摘要。}
本研究文件以活動日捷運需求為背景，整理資料時間界線、正則化迴歸、加法平滑、
樹提升、非負融合及有限期批次候車的數學基礎。對 Ridge 與具有全係數懲罰的
GAM 證明固定設計下解的存在與唯一性；對不限制總和的非負融合，
利用有限生成錐及預測空間的緊緻次水平集證明最佳解存在，
並釐清權重總和改為 2 或 3 的尺度效果。
列車調度以官方班表為固定基準，聯結分鐘需求、剩餘容量、候車積分及增班成本，
給出完整限制、混合整數線性表示及有限候選集合上的存在性證明。
本文區分標準理論、本案推導、既有回顧結果及尚待實作設計；
數學最適性不被解讀為未來預測準確性或實際行車安全保證。

\noindent\textbf{關鍵詞：} 非負融合；強制性；P-spline；批次服務；
混合整數規劃；最佳解存在性。
\section{研究範圍與理論保證}

\subsection{本案要解決的問題}
第一階段利用事前可取得的活動、車站與歷史人流資訊，預測活動日的逐時進站量。
第二階段把固定的需求預測分配為方向與分鐘到達情境，評估列車剩餘容量下的候車，
再搜尋兼顧新增營運成本與候車時間的增班方案。
預測與決策必須分開估計：不能在調度求解時，讓預測權重任意降低需求以節省成本。

本文以目前程式、設定與公開結果為準。既有資料為 21 場活動、四站、每場
18--23 時六個來源時段，共 504 列；主要回顧評估為 2024--2026 年九場北門資料，
共 54 列。同一活動的車站與時段彼此相關，504 列不能視為 504 場獨立活動。
四站為北門、大橋頭、雙連及民權西路。

\begin{table}[htbp]
\centering\small
\caption{方法的實作狀態與本文能證明的內容}
\begin{tabularx}{\linewidth}{@{}p{24mm}p{36mm}X@{}}
\toprule
方法 & 本案狀態 & 理論保證的範圍\\
\midrule
Ridge & 保留作比較與解釋 & 固定設計、正懲罰下，係數解存在且唯一\\
GAM & 保留作非線性比較 & 固定基底、全係數正懲罰下，解存在且唯一\\
XGBoost & 散場加權版本為保留模型 & 固定樹結構的葉節點值可求最適；不保證貪婪建樹全域最優\\
融合 & 曾比較但未採用；新限制為理論設計 & 指定有限資料與損失下的權重解存在性\\
候車情境 & 已有批次上車模擬 & 人數守恆、候車面積與有限期計算\\
最佳化班表 & 公式與教學驗算已完成，實際班表尚未求解 & 有限、非空可行集合中的最佳方案存在\\
\bottomrule
\end{tabularx}
\end{table}

\Needspace{12\baselineskip}
\subsection{「有最佳解」不等於「一定預測準」}
\begin{definition}[四種不同的敘述]
給定問題 $\inf_{x\in\mathcal F}f(x)$：
\begin{enumerate}
\item \textbf{存在性：}存在 $x^*\in\mathcal F$，使 $f(x^*)=\inf_{\mathcal F}f$。
\item \textbf{唯一性：}只有一個決策達到該最小值。
\item \textbf{演算法最適性：}所用演算法確實找到全域最小值，而非僅可行或局部解。
\item \textbf{外部有效性：}對未來活動仍準確，且決策符合未納入模型的營運條件。
\end{enumerate}
\end{definition}

凸性本身不保證最佳解存在。例如 $\min_{x\ge0}e^{-x}$ 是凸問題，
下界為 0，卻沒有任何有限 $x$ 可達到。後文會逐一說明保證存在所需的附加條件。
本文的定理與證明是依標準數學結果針對本案重述或推導，
不是聲稱提出新的普遍性定理。文獻與程式位置列於文末。

\subsection{主要記號}
\begin{table}[htbp]
\centering\small
\begin{tabularx}{\linewidth}{@{}p{23mm}X p{34mm}@{}}
\toprule
記號 & 意義 & 單位或範圍\\
\midrule
$i=(e,s,h)$ & 活動、車站、來源時段的資料列 & $i=1,\ldots,n$\\
$y_i,z_i$ & 真實進站量、前處理後輸入 & 人次；設計向量\\
$\omega_i$ & XGBoost 訓練列權重 & 正數，場內平均為 1\\
$a_i$ & 融合損失的列權重 & $a_i>0,\ \sum_i a_i=1$\\
$w_m$ & 第 $m$ 個模型的融合係數 & 與 $\omega_i$ 不同\\
$x_k$ & 是否選擇第 $k$ 個候選增班 & $\{0,1\}$\\
$Q_d(t),W$ & 方向排隊量、總候車面積 & 人；人分鐘\\
$u,v$ & 單位增量營運成本、時間價值 & 元／列車公里；元／人分鐘\\
\bottomrule
\end{tabularx}
\end{table}

\section{資料時間界線、損失與驗證}

\subsection{資訊可得性與前處理}
令 $\mathcal I_{t_0}$ 為預測截止時間 $t_0$ 已知資訊生成的 $\sigma$-代數；
資料欄位除了觀測時刻，還需記錄真正發布或可取得的時刻。正式日前預測要求
\begin{equation}
\operatorname{available\_at}(z_{ij})\le t_0.
\end{equation}
活動日前一天的實際氣象、事後全月聲量或當日未來時段人流不滿足這個條件。
歷史觀測日期早於 $t_0$，也不代表該資料當時已發布。
本案歷史 OD 與日曆的部分發布時點仍未完全核實，因此目前結果是回顧比較，
不能宣稱已完整驗證嚴格 D$-1$ 部署。

\begin{proposition}[預測資訊的封閉性]
若訓練資料、前處理參數、超參數及目標輸入均可由 $\mathcal I_{t_0}$ 取得，
且計算程序為可測函數，則最終預測亦為 $\mathcal I_{t_0}$ 可測。
\end{proposition}
\begin{proof}
把補值、標準化、編碼、擬合及預測寫成有限個可測映射的合成。
可測映射的合成仍可測；固定亂數種子或將已抽取的種子納入已知資訊即可。
\end{proof}
這只證明沒有把未來資訊放進計算，並不證明預測無偏或泛化誤差小。
本案按年度向前驗證；前處理只由各折訓練資料估計，
2024、2025、2026 年分別使用過去 12、15、18 場活動。
向前驗證的設計參考文獻 \cite{fpp}；九場反覆參與開發的限制仍須保留
\cite{cawley}。

連續欄以訓練平均與母體標準差轉換：
\[
z_{ij}=\frac{x_{ij}-\bar x_{j,\mathrm{train}}}{s_{j,\mathrm{train}}}.
\]
類別與二元欄維持參考編碼。移除常數與完全線性相依欄是識別性處理，
不是依相關係數或 $p$ 值自動選變數。
車站與時段交互作用使時段差異可以因站而異：
\[
\eta_{s,h}=\beta_0+\alpha_s+\gamma_h+\delta_{s,h}+z^\top\beta.
\]
在固定車站 $s$，由 $h_1$ 改為 $h_2$ 的條件差異為
$(\gamma_{h_2}-\gamma_{h_1})+(\delta_{s,h_2}-\delta_{s,h_1})$；
不能把它解讀成活動的因果效果。

\subsection{MSE、MAE 與低估雙倍的不同目標}
定義 $r=y-\hat y$，故 $r>0$ 表示低估。三種損失為
\begin{align}
\ell_2(r)&=r^2,&
\ell_1(r)&=|r|,&
\ell_\kappa(r)&=\kappa[r]_++[-r]_+,\quad\kappa>0.
\label{eq:loss}
\end{align}
本案曾以 $\kappa=2$ 選參；目前 XGBoost 訓練仍使用加權平方誤差，
不能將這兩件事混為「模型已改成分位數訓練」。

\begin{proposition}[損失決定所估的母體位置]
給定條件資訊 $X$，若所需動差有限：
平方誤差的最佳常數預測為 $\E(Y\mid X)$；
絕對誤差的最佳預測為條件中位數；
非對稱絕對誤差的最佳預測為
$\tau=\kappa/(1+\kappa)$ 的條件分位數。
\end{proposition}
\begin{proof}
對平方誤差，設 $\mu=\E(Y\mid X)$，則
\[
\E[(Y-a)^2\mid X]=\Var(Y\mid X)+(a-\mu)^2,
\]
右式在 $a=\mu$ 最小。對非對稱絕對誤差，
若條件分布函數 $F$ 在 $a$ 連續，對 $a$ 微分得
\[
\frac{d}{da}\E\{\ell_\kappa(Y-a)\mid X\}
=F(a)-\kappa\{1-F(a)\}.
\]
零點滿足 $F(a)=\kappa/(1+\kappa)$。
有原子時改用次梯度，條件為
$F(a^-)\le\tau\le F(a)$。令 $\kappa=1$ 即得中位數結論。
\end{proof}
因此低估罰 2 倍的母體目標是 $2/3$ 分位數，而非自動等於平均人流。
這是條件母體目標的敘述；受限的三模型融合不一定能精確表示該分位數。

回顧評分採
\[
\operatorname{MAE}=\frac1N\sum_i|\hat y_i-y_i|,\qquad
\operatorname{RMSE}=\sqrt{\frac1N\sum_i(\hat y_i-y_i)^2}.
\]
RMSE 與 MSE 在同一資料、相同權重下的排序相同，因平方根嚴格遞增。
這不表示 RMSE 與 MAE 的最小點相同，也不表示使用平方誤差訓練一定
能最佳化實際列車調度成本。

\subsection{有限調參與選模的界線}
固定有限候選超參數集合 $\Theta$，若每個候選皆成功產生有限驗證損失，
則 $\argmin_{\theta\in\Theta}\widehat R_{\mathrm{val}}(\theta)$ 非空。
此證明只需在有限個實數中取最小值，不能推論找到連續參數空間的最適解。
失敗候選需明確排除並記錄；若全部失敗，便無可比較的模型。
反覆觀看同一批驗證結果會帶來選模偏差，故「存在最小驗證誤差」
不等於「已無偏估計未來誤差」\cite{cawley}。

\section{Ridge：線性結構、收縮與唯一解}

\subsection{目標與閉式解}
令 $Z\in\R^{n\times p}$ 為固定訓練設計，截距不受懲罰。
本案採用
\begin{equation}
\min_{\beta_0,\beta}
\frac1n\norm{y-\beta_0\one-Z\beta}_2^2+\lambda\norm{\beta}_2^2,
\qquad\lambda>0.
\label{eq:ridge}
\end{equation}
這是 Ridge 的經典設定 \cite{ridge}，其懲罰作用是穩定係數，
不保證每一個係數的絕對值都相對 OLS 變小，也不保證所有資料的預測誤差都下降。

\begin{theorem}[Ridge 解存在且唯一]
若 $n\ge1$、$Z,y$ 有限且 $\lambda>0$，式 \eqref{eq:ridge}
有唯一最佳解，無須假設 $Z$ 的欄向量線性獨立。
\end{theorem}
\begin{proof}
對固定 $\beta$，截距的一階條件給出
$\hat\beta_0=\bar y-\bar z^\top\beta$。
令 $Z_c=Z-\one\bar z^\top$、$y_c=y-\bar y\one$，
問題等價於最小化
$n^{-1}\norm{y_c-Z_c\beta}^2+\lambda\norm{\beta}^2$。
對任意非零 $v$，
\[
v^\top(Z_c^\top Z_c+n\lambda I)v
=\norm{Z_cv}^2+n\lambda\norm v^2>0.
\]
故係數矩陣正定，唯一駐點為
\begin{equation}
\hat\beta=(Z_c^\top Z_c+n\lambda I)^{-1}Z_c^\top y_c,\qquad
\hat\beta_0=\bar y-\bar z^\top\hat\beta .
\end{equation}
目標為嚴格凸二次函數，該駐點即唯一全域最小點。
\end{proof}

若程式採用未除以 $n$ 的 SSE，
必須用 $\alpha_{\mathrm{software}}=n\lambda$，才能與
\eqref{eq:ridge} 相同。本案用 SVD 求解，不直接以數值反矩陣實作。

\subsection{收縮的幾何意義}
令 $Z_c=U\operatorname{diag}(d_j)V^\top$ 為 SVD。
在有資料支持的奇異方向上，
\[
Z_c\hat\beta
=U\operatorname{diag}\!\left(\frac{d_j^2}{d_j^2+n\lambda}\right)U^\top y_c .
\]
每個倍率介於 0 與 1；小奇異值方向受到較強收縮，
所以近共線性引起的不穩定可以被限制。
這是在正交方向上的收縮，不是每個原始變數係數的逐項單調性定理。

最終人流輸出採 $\hat y=\max(0,\eta)$。
若真值 $y\ge0$，對任意 $\eta<0$，有
$|y-0|=y\le y-\eta=|y-\eta|$，
所以截去負預測不會增加該筆絕對或平方誤差；
這不等於在融合求解時任意加入截斷後仍保有原來的凸性。

\section{GAM：有限基底、平滑懲罰與識別性}

\subsection{本案使用的加法模型}
Gaussian identity GAM 允許部分連續變數使用加法曲線：
\begin{equation}
\eta_i=\beta_0+x_i^\top\beta+\sum_{j\in\mathcal S}f_j(z_{ij}),
\qquad f_j(z)=\sum_{r=1}^{K_j}\theta_{jr}B_{jr}(z).
\end{equation}
將線性項與樣條基底合併為矩陣 $A$，所有非截距係數合併為 $\theta$。
本案使用 cubic P-spline；參考 GAM 與 P-spline 文獻
\cite{gam,pspline}，但支持門檻及格點是自行設定。
連續候選須至少六場已知資料；指定可平滑欄另須至少八個相異已知值。
一般日基準在現行設定為線性項，並非所有連續欄都畫成曲線。

令 $D_2$ 為二階差分矩陣，樣條係數的差分為
$\theta_{r+2}-2\theta_{r+1}+\theta_r$。目標為
\begin{equation}
\min_{\beta_0,\theta}
\norm{y-\beta_0\one-A\theta}^2
+\lambda_s\sum_{j\in\mathcal S}\norm{D_2\theta_j}^2
+n\alpha_G\norm{\theta}^2,
\label{eq:gam}
\end{equation}
其中 $\lambda_s\ge0,\alpha_G>0$。係數差分懲罰不是
精確的 $\int[f_j''(z)]^2\,dz$；兩者不可混寫為同一個實作。

\begin{theorem}[目前 GAM 的唯一解]
在固定有限基底、有限資料及 $\alpha_G>0$ 下，
式 \eqref{eq:gam} 對截距與所有非截距係數有唯一最佳解。
\end{theorem}
\begin{proof}
將差分懲罰組成半正定矩陣 $S$，使
$\theta^\top S\theta=\lambda_s\sum_j\norm{D_2\theta_j}^2$。
對任意變動 $(b,v)\ne(0,0)$，目標 Hessian 的二分之一二次型為
\[
\norm{b\one+Av}^2+v^\top Sv+n\alpha_G\norm v^2 .
\]
若 $v\ne0$，最後一項嚴格正；若 $v=0,b\ne0$，
第一項為 $nb^2>0$。故 Hessian 正定，
二次函數具有唯一全域最小點。
\end{proof}

全係數 L2 懲罰是上述證明的重要條件。
只靠二階差分懲罰時，常數與線性趨勢落在其零空間，
若再與截距或其他項重疊，不能直接宣稱係數唯一。
實際程式另核對懲罰矩陣、收斂狀態與正規方程殘差。

\subsection{解釋與外插的限制}
平滑分量的中心化只是圖形呈現慣例，
不能把單條曲線當成控制所有混雜後的因果效果。
有限基底數與正懲罰能保證固定問題的唯一解，
不保證邊界外的曲線合理；本案邊界外仍可能線性外插。
「GAM 融合權重為 0」也不是此定理失效，
因為基礎模型擬合與後續融合權重是兩個不同最佳化問題。

\section{XGBoost：加權樹提升與能證明的局部步驟}

\subsection{散場訓練權重}
對活動 $e$ 的資料列，令 $o_i\in[0,1]$ 為事前預定結束後兩小時與該來源時段的重疊，
定義
\begin{equation}
r_i=1+3\,\mathbf1(s_i=\text{北門})o_i,\qquad
\omega_i=\frac{n_e r_i}{\sum_{j\in e}r_j}.
\label{eq:weights}
\end{equation}
因此 $\sum_{i\in e}\omega_i=n_e$，每場平均權重為 1，
避免只因提高北門權重而使整場總權重任意放大。
未知預定結束時間不加權。
這個「最多四倍的原始列權重」不等於損失中「低估罰兩倍」，
亦不等於三模型融合權重。

\subsection{固定樹結構下的葉節點最適值}
在第 $t$ 次提升中，令舊預測為 $\hat y_i^{(t-1)}$、
新增樹為 $f_t$。以加權半平方誤差與樹懲罰表示：
\begin{equation}
\sum_i\frac{\omega_i}{2}
\{y_i-\hat y_i^{(t-1)}-f_t(z_i)\}^2
+\gamma_T J+\frac{\lambda_T}{2}\sum_{j=1}^J v_j^2.
\end{equation}
$J$ 是葉節點數、$v_j$ 是葉節點值；這裡的 $\lambda_T$
不是 Ridge 的 $\lambda$。依 XGBoost 原文第 2 節的二階表示
\cite{xgb}，設
\[
g_i=\omega_i(\hat y_i^{(t-1)}-y_i),\quad h_i=\omega_i,\quad
G_j=\sum_{i\in I_j}g_i,\quad H_j=\sum_{i\in I_j}h_i .
\]
對平方損失，此二階展開除常數外是精確等式，而不只是局部近似。
固定各葉的樣本集合 $I_j$ 後，目標分離為
\[
\sum_{j=1}^J\left[G_jv_j+\frac12(H_j+\lambda_T)v_j^2\right]+\gamma_TJ .
\]

\begin{proposition}[固定分區的唯一葉值]
若每葉 $H_j+\lambda_T>0$，則
\begin{equation}
v_j^*=-\frac{G_j}{H_j+\lambda_T}
\label{eq:leaf}
\end{equation}
是該葉唯一最佳值。
\end{proposition}
\begin{proof}
對 $v_j$ 微分得 $G_j+(H_j+\lambda_T)v_j$；
二階導數嚴格正，故令一階導數為零即得唯一全域最小點。
各葉可分離求解，得到固定樹結構的最佳葉值。
\end{proof}

候選分裂的目標改善量為
\begin{equation}
\operatorname{Gain}
=\frac12\left[
\frac{G_L^2}{H_L+\lambda_T}
+\frac{G_R^2}{H_R+\lambda_T}
-\frac{(G_L+G_R)^2}{H_L+H_R+\lambda_T}\right]-\gamma_T .
\end{equation}
現行模型的 L1 懲罰為 0，故採上述葉值式；
若另加 L1，分子須改為對 $G_j$ 的軟閾值。
學習率 $\nu$ 再縮小新增樹：
$\hat y^{(t)}=\hat y^{(t-1)}+\nu f_t$。

\subsection{不能由葉值定理推出整個建樹全域最優}
貪婪建樹逐次選分裂，hist 限制可考慮的切點，
並受樹深、葉節點門檻與樹數限制。
因此式 \eqref{eq:leaf} 不保證找出所有樹結構與所有提升步驟的聯合全域最佳解。
若列舉有限個樹結構，每個結構有最小值，可證明這個有限集合內至少有最佳結構；
但實際 XGBoost 並未窮舉該集合。
同理，調參的有限格點最小值存在，不代表已搜尋所有可能設定。
目前保留版本沒有採用 early stopping。

\section{融合權重：總和自由、總和為 2 或 3}

\subsection{固定留出預測上的融合問題}
把 Ridge、GAM、XGBoost 在同一批過去留出資料上的預測合為
$P=[p_R,p_G,p_X]\in\R^{n\times3}$。此處基礎模型及其預測固定，
只估計融合係數：
\begin{equation}
\min_{w\ge0}F(w)
=\sum_{i=1}^n a_i\ell\{y_i-(Pw)_i\},
\qquad a_i>0,\quad\sum_i a_i=1 .
\label{eq:fusion}
\end{equation}
若活動等權，每場有 $n_e$ 列、共 $N_E$ 場，
可令該場每列 $a_i=1/(N_En_e)$。
這些 $a_i$ 與式 \eqref{eq:weights} 的訓練權重不同。
非負且不限制總和的融合有文獻先例
\cite{breiman,pystacked}；總和為 1 是一種模型限制，不是所有 stacking 的公理。

\begin{lemma}[有限生成錐封閉]\label{lem:cone}
對有限矩陣 $P$，集合
$C=\{Pw:w\ge0\}=\operatorname{cone}\{p_R,p_G,p_X\}$
是非空封閉凸集合。
\end{lemma}
有限生成錐的標準結果見 \cite{bertsekas}；
本案所需的完整證明另列附錄 \ref{app:cone}。

\begin{theorem}[總和自由時的存在性]\label{thm:free}
若 $P,y$ 有限，$\ell$ 取平方、絕對或式 \eqref{eq:loss}
中 $\kappa>0$ 的非對稱絕對損失，
則式 \eqref{eq:fusion} 至少有一組全域最佳權重，
不要求 $P$ 的三個欄向量線性獨立。
\end{theorem}
\begin{proof}
改在可達預測集合 $C$ 上最小化
$g(z)=\sum_i a_i\ell(y_i-z_i)$。
令 $a_{\min}=\min_i a_i>0$。三種損失依序滿足
\[
g(z)\ge a_{\min}\norm{y-z}_2^2,\qquad
g(z)\ge a_{\min}\norm{y-z}_1,\qquad
g(z)\ge a_{\min}\min(1,\kappa)\norm{y-z}_1 .
\]
所以 $g$ 連續，且 $\norm z\to\infty$ 時 $g(z)\to\infty$。
令
\[
K=\{z\in C:g(z)\le g(0)\}.
\]
由 $0\in C$，$K$ 非空；由引理 \ref{lem:cone} 與連續性，$K$ 封閉；
由上述不等式，$K$ 有界。有限維空間中 $K$ 緊緻，
Weierstrass 極值定理保證 $g$ 在 $K$ 取得最小值 $g(z^*)$。
$K$ 外的點比可行點 0 更差，故 $z^*$ 亦為 $C$ 上最小點。
由 $z^*\in C$，存在有限 $w^*\ge0$ 使 $Pw^*=z^*$，
即得到式 \eqref{eq:fusion} 的最佳解。
\end{proof}

這個證明在\textbf{預測空間}建立有界次水平集；
沒有誤稱權重集合 $\R_+^3$ 本身有界。
它也不要求資料彼此獨立；資料相依影響泛化評估與統計推論，
但不會使固定有限資料上的這個存在性定理失效。

\subsection{唯一性、KKT 與線性規劃}
對平方損失，令 $D_a=\diag(a_i)$，
\[
\nabla F(w)=2P^\top D_a(Pw-y),\qquad
\nabla^2F(w)=2P^\top D_aP\succeq0 .
\]
若 $\rank P=3$，權重解唯一；否則仍存在最小點，但可能有多組權重。
因平方損失在預測向量上嚴格凸，最適訓練擬合 $Pw^*$ 唯一，
不代表各組解在新輸入上仍產生相同預測。
必要且充分的 KKT 條件為
\begin{equation}
w^*\ge0,\quad g^*=2P^\top D_a(Pw^*-y)\ge0,\quad
w_m^*g_m^*=0 .
\end{equation}
此為凸二次問題的最適性條件 \cite{boyd}。
必要性可由可行方向看出：$w_m^*>0$ 時可向兩側微動，故 $g_m^*=0$；
$w_m^*=0$ 時可向正方向微動，故 $g_m^*\ge0$。
充分性由凸函數的一階不等式得到：對任意 $w\ge0$，
\[
F(w)\ge F(w^*)+(g^*)^\top(w-w^*)\ge F(w^*),
\]
因互補條件使 $(g^*)^\top w^*=0$，且 $(g^*)^\top w\ge0$。

非對稱絕對損失可改寫為線性規劃：
\begin{equation}
\begin{aligned}
\min_{w,r^+,r^-}\quad&\sum_i a_i(\kappa r_i^++r_i^-),\\
\text{使得}\quad&y-Pw=r^+-r^-,\\
&w,r^+,r^-\ge0 .
\end{aligned}
\end{equation}
$\kappa=1$ 為 MAE，$\kappa=2$ 為低估雙倍。
不可微不妨礙全域最佳化。
對固定 $w$，若同一列 $r_i^+,r_i^-$ 同時為正，
兩者同減可維持等式並嚴格降低成本。
故最小時分別等於殘差的正、負部分，證明上述 LP 與原損失等價。
若另外加入 $\eta\norm w^2$、$\eta>0$，
目標在權重空間強制且嚴格凸，因此三種損失皆有唯一權重解；
這是新的融合層正則化，仍須用過去資料選擇 $\eta$。

\subsection{權重總和等於 2 或 3}
\begin{theorem}[固定正總和的存在性]\label{thm:sum}
對任意固定 $c>0$，
\[
\min_{w\in S_c}F(w),\qquad
S_c=\{w\ge0:\one^\top w=c\}
\]
若 $F$ 連續，則至少有一組最佳解。特別是 $c=2$ 或 $3$ 均成立。
\end{theorem}
\begin{proof}
$(c,0,0)\in S_c$，故非空。
等式與非嚴格不等式定義的 $S_c$ 封閉，
且每個分量滿足 $0\le w_m\le c$，故有界。
Weierstrass 極值定理立即給出最小點。
\end{proof}

\textbf{重新最佳化後的比例未必不變。}
例如教學資料 $n=1$，真值為 1，三個預測為 $(1,2,0.2)$。
總和為 1 時 $(1,0,0)$ 可達零誤差；
把它乘 2 得 $(2,0,0)$，在總和為 2 的問題中誤差為 1，
但 $(0.75,0,1.25)$ 的預測仍為 1，可達零誤差。
因此「把原權重乘 2」不一定是新問題的最佳解；
存在性與最佳比例是否相同是不同問題。

若總和固定為 $c$，卻用 $Pw/c$ 計算訓練損失，
令 $\alpha=w/c$ 就有 $\alpha\ge0,\sum_m\alpha_m=1$，
所以與總和為 1 的加權平均完全等價。
此等價要求目標同步轉換；若原來加上 $\eta\norm w^2$，
變數替換後是 $\eta c^2\norm\alpha^2$，不能忽略懲罰尺度。

若三模型皆預測 1,000 人，直接使用總和為 2 或 3 的 $Pw$，
必分別預測 2,000 或 3,000 人。數學上仍有最小點，
但限制會強制改變整體尺度。總和自由時，
只要 $s=\sum_mw_m>0$，可寫成
\[
Pw=sP\alpha,\qquad \alpha=w/s,\quad\sum_m\alpha_m=1,
\]
表示同時學習配比與尺度；$s=0$ 為零預測的特例。

\subsection{嚴格正權重、固定數值與誤差互補}
要求 $w_m>0$ 時集合不再封閉，存在性可能失敗。
例如 $P=(1,1,1)$、$y=0$，平方誤差為
$(w_R+w_G+w_X)^2$；三個權重嚴格正時永遠大於 0，
但可一起趨近 0。因此有下界卻沒有最小點。
若改成 $w_m\ge\varepsilon>0$，
令 $w=\varepsilon\one+u$，即可在總和自由時套用定理 \ref{thm:free}；
總和固定為 $c$ 時則須先有 $3\varepsilon\le c$。

若全部權重固定為 $(2,3,4)$，只有一個候選，沒有要學習的配比；
若只固定 $w_R=2$，其餘非負且總和自由，
把目標改為 $y-2p_R$ 即可對剩餘兩欄套用同一存在性證明。
每個權重從 $\{2,3\}$ 選最多 8 組，有限集合內亦有最佳組合，
但不等於連續問題的最佳解。

令 $e_m=p_m-y$，誤差矩陣為 $E=[e_R,e_G,e_X]$。
當 $\sum_mw_m=1$ 時，
\begin{equation}
Pw-y=Ew,\qquad
\operatorname{MSE}(w)=w^\top(E^\top E/n)w .
\end{equation}
$E^\top E/n$ 是未中心化二階矩，不只是相關矩陣。
設誤差均值向量為 $\mu$、以 $1/n$ 定義的樣本共變異矩陣為 $\Sigma$，
則 $E^\top E/n=\Sigma+\mu\mu^\top$。
因此高誤差相關與共同偏差會限制融合的互補效益。
若權重總和為 $s\ne1$，則
\[
Pw-y=Ew+(s-1)y,
\]
不能繼續沿用總和為 1 的純誤差共變異式。
最小化過去資料損失也不保證未來優於 XGBoost，
故現行專案沒有因存在性定理就改採融合。

\section{從逐時預測到有限期批次候車}

\subsection{需求輸入與分鐘情境}
對北門第 $h$ 小時，令 $\widehat Y_h\ge0$ 為已凍結的人流預測。
目前保留模型是散場加權 XGBoost；若日後採用融合，則用
$\widehat Y_h=\sum_m\widehat w_m\widehat y_{hm}$，但權重必須先在預測階段估計。
設方向比例 $\pi_{dh}\ge0$、小時內分鐘比例 $p_{hr}\ge0$，且
$\sum_d\pi_{dh}=1$、$\sum_{r=1}^{60}p_{hr}=1$，定義
\begin{equation}
A_{dhr}=\widehat Y_h\pi_{dh}p_{hr},\qquad
\sum_{d,r}A_{dhr}=\widehat Y_h.
\label{eq:arrival}
\end{equation}
均勻情境取 $p_{hr}=1/60$；寬峰與集中峰取不同的非負形狀再正規化。
這些形狀是對未知分鐘分布的情境假設，不是由逐時總量唯一識別的結果。
無窮多組分鐘曲線可能產生相同逐時總量。
若另納入閘門至月台步行延遲，須使用總質量為 1 的延遲分布，
並保留跨出小時或模擬終點的尾端人流。

\textbf{兩階段不能混成任意降低需求的最佳化。}
若把融合權重 $w\ge0$ 和列車決策一起放入只有營運、候車成本的目標，
可能藉由 $w=0$ 得到零需求，製造虛假的低成本。
本文先估計並固定需求，再比較相同需求下的班表。

\subsection{剩餘容量與事件格點}
列車 $j$ 的可上車容量為
\begin{equation}
R_j=\max\{0,K_{\mathrm{use},j}-(O_j-H_j)\},\qquad 0\le H_j\le O_j,
\end{equation}
其中 $O_j$ 是到站前車上人數、$H_j$ 是本站下車人數，
$K_{\mathrm{use},j}$ 是採用的載客上限。
六節高運量列車的最大參考容量不等於每班到北門後都還能載入同樣人數。
目前逐班剩餘容量尚未實測，555、1,110、1,665 人為本案情境值。
新增班次也不自動假設抵站時為空車。

對方向 $d$，合併分鐘到達率轉折、原班列車、所有候選列車與共同終點 $T$，
建立固定格點
\[
0=t_{d,0}<t_{d,1}<\cdots<t_{d,J_d}=T,\qquad
\Delta_{d,j}=t_{d,j}-t_{d,j-1}.
\]
$q_{d,0}$ 是起點服務處理完後的排隊量，
$a_{d,j}$ 是區間 $(t_{d,j-1},t_{d,j}]$ 的到達量；該區間內到達率固定。
令 $B_d$ 為方向 $d$ 的原班表，$E$ 為有限個候選增班，
$\tau_k$ 為到達北門時刻，$x_k\in\{0,1\}$ 表示是否選用。
格點服務容量為
\begin{equation}
C_{d,j}(x)=C^0_{d,j}
+\sum_{k\in E:\,d_k=d,\,\tau_k=t_{d,j}}R_kx_k,
\quad
C^0_{d,j}=\sum_{b\in B_d:\,\tau_b=t_{d,j}}R_b.
\label{eq:capacity}
\end{equation}
乘客可在服務時刻前持續累積；列車到站後以批次上車處理。

\subsection{人數守恆與總候車時間}
\begin{equation}
\boxed{\begin{aligned}
q^-_{d,j}&=q_{d,j-1}+a_{d,j},\\
b_{d,j}&=\min\{q^-_{d,j},C_{d,j}(x)\},\\
q_{d,j}&=q^-_{d,j}-b_{d,j}
=\max\{0,q_{d,j-1}+a_{d,j}-C_{d,j}(x)\}.
\end{aligned}}
\label{eq:queue}
\end{equation}
這是有限期流體需求與批次服務模型；有人排隊且車上有剩餘容量時就上車。
沒有上滿可能來自該方向當時候車量不足；並非模型要求列車必須留空。

\begin{proposition}[人數守恆]
對任一方向與任一 $j$，
\[
q_{d,j}+\sum_{\ell=1}^j b_{d,\ell}
=q_{d,0}+\sum_{\ell=1}^j a_{d,\ell}.
\]
且 $q_{d,j},b_{d,j}\ge0$，$b_{d,j}\le C_{d,j}(x)$。
\end{proposition}
\begin{proof}
由式 \eqref{eq:queue} 有 $q_{d,j}-q_{d,j-1}=a_{d,j}-b_{d,j}$；
對 $j$ 次遞推求和即得。非負性與容量限制直接由 min 與 max 的定義成立。
\end{proof}

區間內沒有其他服務，排隊量從 $q_{d,j-1}$ 線性增加，故
\begin{equation}
\boxed{W(x)=\sum_d\int_0^T Q_d(t;x)\,dt
=\sum_{d,j}\left[q_{d,j-1}\Delta_{d,j}
+\frac12 a_{d,j}\Delta_{d,j}\right].}
\label{eq:waiting}
\end{equation}
$W$ 的單位為人分鐘。只加總列車離站後排隊量會漏掉班次之間的等待，
亦缺少時間單位。

\begin{theorem}[候車面積的樣本路徑恆等式]
假設起點無排隊，旅客 $\ell$ 的到達與上車時刻分別為
$A_\ell\in[0,T]$、$B_\ell\ge A_\ell$。
未上車者令 $B_\ell=+\infty$。對有限個旅客，
\begin{equation}
\int_0^T Q(t)\,dt
=\sum_\ell\{\min(B_\ell,T)-A_\ell\}.
\label{eq:area}
\end{equation}
若全部在 $T$ 前上車，右側就是總完整候車時間。
\end{theorem}
\begin{proof}
除不影響積分的端點外，
$Q(t)=\sum_\ell\mathbf1\{A_\ell\le t<B_\ell\}$。
交換有限求和與積分，每個指示函數的積分是
$[A_\ell,B_\ell)$ 與 $[0,T]$ 的交集長度，得到式 \eqref{eq:area}。
對流體需求，將有限求和改成非負質量積分，Tonelli 定理給出相同結論。
\end{proof}

這個面積恆等式不需要 Poisson 到達、指數服務或穩態；
與 Little 定律的計數思路相關，但不能把有限期間截尾量
直接當作無限期間穩態平均 \cite{little,urban}。
若起點已有旅客，面積只含他們在起點之後的等待。
若起點為零且期末全員上車，平均候車為
$W/\sum_{d,j}a_{d,j}$（分母須大於零）；否則應另報期末滯留，
不能宣稱已估計完整平均等待。

\subsection{為何不直接使用 M/M/c}
M/M/c 的 $c$ 是同時並行、逐人服務的服務台數，
並非一小時到站的列車班數。其穩態結果通常要求
固定 Poisson 到達率 $\lambda$、獨立指數服務率 $\mu$，且
$\rho=\lambda/(c\mu)<1$。
本案散場需求非定態、列車定時批次提供容量，所以採式 \eqref{eq:queue}。
分鐘分配只是需求建模的一步；排隊模型還包含容量、上車、滯留與等待時間。
短時間需求高於供給時，仍能計算有限期滯留，
但不能因此聲稱對應的 M/M/c 穩態等待有限。

\section{班表最佳化：完整問題、存在性與求解形式}

\subsection{基準班表、提前量與班距}
原班表固定使用已整理的北門假日時刻表；這是官方發布時的研究基準，
不等於已取得每個歷史活動日的實際行車紀錄。
新增候選包含方向、到站時刻、起站至北門的行駛時間 $\delta_k$、
整備時間 $\ell_k$、完整新增公里 $L_k$ 及剩餘容量 $R_k$。
若決策時刻是 $t_0$，候選至少滿足
\begin{equation}
r_k=\tau_k-\delta_k-\ell_k\ge t_0.
\end{equation}
若尚無整備資料而暫取 $\ell_k=0$，這只是較樂觀的研究假設。
起站到北門的時間是提前量的一部分，不是完整增班行駛里程。

令研究最小同方向班距為 $h_{\min}$，定義 $e_k\in\{0,1\}$：
只有提前量合格，且與該方向所有相關原班時刻相隔至少 $h_{\min}$，才令 $e_k=1$。
相鄰視野外原班也要納入衝突檢查。
令 $\mathcal P$ 為同方向且 $|\tau_k-\tau_l|<h_{\min}$ 的候選對，則
\begin{equation}
\mathcal X=\left\{x\in\{0,1\}^{K}:
x_k\le e_k,\quad x_k+x_l\le1\ ((k,l)\in\mathcal P),\quad
x_k=\bar x_k\ (k\in\mathcal C)\right\}.
\label{eq:feasible}
\end{equation}
$\mathcal C$ 是已承諾、不可取消的候選，$\bar x_k$ 為其固定決策。
同時刻的互斥候選也放入 $\mathcal P$。
若另有跨方向共用路段或折返限制，須再增加限制，不能由此式自動推得安全。

目前以 3 分鐘作研究門檻，來源是官方尖峰平均班距資訊；
\textbf{平均班距不能當成已核實的最小安全班距}。
備車數未取得，依目前研究範圍未設定備車供給限制；
因此求得方案仍可能超出實際車隊、司機或折返能力。
官方來源索引見附錄 \ref{app:implementation}。

\subsection{營運成本與時間成本的聯合目標}
新增班次成本先寫為 $c_k=uL_k$，$u$ 單位為元／列車公里。
$L_k$ 應含完整新增服務與需計價的回送，不只計算到北門為止。
目前 $u=128.05$ 是使用者提供的研究試算值，
成本與公里的年度、車公里或列車公里口徑尚須統一核實。
對所有方案都不變的年度固定成本，可從目標刪去；
只有會隨增班改變的啟用成本，才另列入 $c_k$。

時間價值取 $v=196/60$ 元／人分鐘，
是以 2026 年最低時薪作代理值 \cite{mol}，不是已估計的乘客願付價格。
可比較 $0.5v,v,2v$，且與營運成本使用相同幣值年度。
候車成本是貨幣化的不便，不是業者實際付出的現金。

以期末排空為要求，完整問題為
\begin{equation}
\boxed{\begin{aligned}
\min_{x,q,b}\quad
&\sum_{k=1}^{K}uL_kx_k
+v\sum_{d,j}\left(q_{d,j-1}\Delta_{d,j}
+\tfrac12 a_{d,j}\Delta_{d,j}\right)\\
\text{使得}\quad
&x\in\mathcal X,\qquad C_{d,j}(x)\ \text{由式 \eqref{eq:capacity} 決定},\\
&q^-_{d,j}=q_{d,j-1}+a_{d,j},\\
&b_{d,j}=\min\{q^-_{d,j},C_{d,j}(x)\},\\
&q_{d,j}=q^-_{d,j}-b_{d,j},\qquad q_{d,J_d}=0\quad\forall d.
\end{aligned}}
\label{eq:dispatch}
\end{equation}
需求 $a_{d,j}$、原班表及比較終點 $T$ 對所有方案相同。
相較原班表的目標變化可寫成
$\sum_k c_kx_k+v\{W(x)-W(0)\}$；
因 $W(0)$ 為固定常數，最適班表與式 \eqref{eq:dispatch} 相同。
新增班數為 $\sum_kx_k$；這個計數差不能取代等待時間和行駛公里的比較。

共同考慮候車與營運成本的研究方向見 Pan 等 \cite{pan}。
本節是針對北門的簡化模型與自行推導，並非該論文的原式；
本次核對其期刊摘要，沒有聲稱重現全文的所有限制。
只計算北門等待時，結論亦只對本研究邊界成立，不能稱整條路線總效益最優。

\subsection{最佳解存在性}
\begin{theorem}[有限候選班表的最小值]
假設候選數 $K$、事件格點數及 $T$ 有限；
起始排隊、到達、容量與成本皆有限且非負；
式 \eqref{eq:dispatch} 至少有一個可行班表。
則至少存在一個全域最佳班表。
\label{thm:dispatch}
\end{theorem}
\begin{proof}
二元向量最多有 $2^K$ 個，滿足全部限制的集合是其非空有限子集。
對每個可行 $x$，式 \eqref{eq:queue} 逐步唯一決定排隊與上車量。
令方向總人數上界為
$D_d=q_{d,0}+\sum_j a_{d,j}<\infty$。
由守恆可得 $0\le Q_d(t;x)\le D_d$，因此
\[
0\le W(x)\le T\sum_dD_d<\infty.
\]
新增成本也因候選數與單班成本有限而有限。
非空有限集合上的有限實值函數必取得最小值，
故存在最佳 $x^*$，再由遞推得到相應 $q^*,b^*$。
\end{proof}

此證明不要求二元問題凸或可微，也不保證最佳解唯一或求解迅速。
它保證候選時間格點內有最佳方案，不保證對所有秒級連續發車時刻最優。
若改成連續時刻，須另證可行集合非空、緊緻及目標的下半連續性。

\textbf{非空前提不能省略。}
總服務容量不足或最後一班後仍有旅客到達時，可能無法排空。
總容量足夠也不充分，因服務可能發生在旅客抵達之前。
無解時應報告是哪個限制造成，不能免費補尾班車或提早截斷時間軸。

若允許期末滯留，必須明確換成另一個問題：
\begin{equation}
\min_{x\in\mathcal X}\left\{
\sum_kc_kx_k+vW(x)+\gamma\sum_d q_{d,J_d}(x)\right\},\qquad \gamma\ge0.
\label{eq:softterminal}
\end{equation}
移除期末排空等式、保留排隊遞推；$\gamma$ 為元／滯留人次。
只要 $\mathcal X$ 非空且 $\gamma$ 有限，同一證明仍成立。
有限懲罰不保證排空，$\gamma=0$ 則只衡量視野內等待。

\subsection{精確的混合整數線性表示}
令 $s_{d,j}=q_{d,j-1}+a_{d,j}-C_{d,j}(x)$，
以及 $\overline C_{d,j}=C^0_{d,j}+\sum_{k:\,d_k=d,\,\tau_k=t_{d,j}}R_k$。
引入 $z_{d,j}\in\{0,1\}$，用以下限制替代 max：
\begin{equation}
\begin{aligned}
0\le q_{d,j}&\le D_d, &q_{d,j}&\ge s_{d,j},\\
q_{d,j}&\le D_dz_{d,j}, &
q_{d,j}&\le s_{d,j}+\overline C_{d,j}(1-z_{d,j}).
\end{aligned}
\label{eq:milp}
\end{equation}
\begin{proposition}[式 \eqref{eq:milp} 的等價性]
對遵守前段守恆的狀態，式 \eqref{eq:milp} 與
$q_{d,j}=\max(0,s_{d,j})$ 完全等價。
\end{proposition}
\begin{proof}
$z=1$ 時強制 $q=s\ge0$；$z=0$ 時強制 $q=0$ 且 $s\le0$。
反之，$s\ge0$ 選 $z=1$，由守恆 $s\le D_d$；
$s\le0$ 選 $z=0$，因 $s\ge-\overline C_{d,j}$，最後一條亦成立。
由起點逐段套用即可。
\end{proof}
所有上界都由需求與容量計算，無須任意設定巨大常數。
式 \eqref{eq:capacity} 與 \eqref{eq:waiting} 為線性，
搭配式 \eqref{eq:feasible}、\eqref{eq:milp} 及期末條件，得到 MILP。
上車量可由 $b=q_{\mathrm{prev}}+a-q$ 還原。
只保留 $q\ge s,q\ge0$ 一般不是服務規則的完整等價限制，
除非另證所用目標會讓相關不等式都取緊。

\subsection{滾動更新與「等於 2、3」的不同意思}
滾動更新時，先固定當下排隊、已發車或承諾的列車及當時可用需求，
求解未來有限視野，僅執行近期決策，再更新狀態重新求解。
這是模型預測控制的設計方向；目前尚未完成實際動態最佳班表。
每次子問題存在最小值，不代表下一次一定可行，
更不保證所有未來資訊已知時的全期間最優。

\begin{table}[htbp]
\centering\small
\begin{tabularx}{\linewidth}{@{}p{35mm}X@{}}
\toprule
設定 & 數學意義與條件\\
\midrule
$v=2$ 或 $v=3$ & 固定時間價值；定理 \ref{thm:dispatch} 仍適用，不必與其他成本係數加總為 1。\\
$\sum_kx_k=2$ 或 $3$ & 恰好增開兩班或三班；集合仍有限，但可能因班距、提前量或排空要求而為空。\\
$\sum_kx_k\le2$ 或 $3$ & 最多增開指定班數；仍須檢查完整可行性。\\
$x_k=2$ 或 $3$ & 違反本案二元候選定義；要增加不同時刻的候選，不能在同一時刻直接乘載客量。\\
\bottomrule
\end{tabularx}
\end{table}
若寫成 $\alpha\Delta C+\beta vW$，事先固定的無量綱偏好係數
$\alpha,\beta>0$ 不須總和為 1；共同乘同一正數不改最適解，
改變兩者比例才改變取捨。不能把它們也當自由非負變數一起最小化，
否則全取 0 就能消除目標。

\section{既有實驗與理論結論的對照}

以下為專案已保存的回顧結果，本文件沒有重新訓練模型或新增北門調度實驗。
同一批 2024--2026 年九場北門活動、每場六個來源時段的比較如下。
\begin{table}[htbp]
\centering
\caption{既有北門回顧誤差；單位皆為人次／來源小時}
\begin{tabular}{@{}lrr@{}}
\toprule
模型 & MAE & RMSE\\
\midrule
Ridge &845.44&1547.49\\
GAM &830.71&1315.75\\
XGBoost（未加散場權重） &473.91&766.90\\
散場加權 XGBoost（目前保留） &440.28&690.91\\
\bottomrule
\end{tabular}
\end{table}

2026 年 Ridge 的 $\lambda=0.177827941004$，訓練列數 $n=432$，
故程式內 $\texttt{alpha}=n\lambda\approx76.82167$。
GAM 各外層年度使用 $\lambda_s=0.1,\alpha_G=0.001$，
基底數在 2024、2025 年為 5、2026 年為 7。
保留的 XGBoost 在 2024 年使用深度 3、250 棵樹、學習率 0.1；
2025、2026 年為深度 4、250 棵樹、學習率 0.05。
這些都是各年度可用過去資料下的設定，不是用九場真值重新一次擬合的單一模型。

這九場已反覆用於開發，故結果支持「目前比較中選擇保留此模型」，
尚不足以證明對全新活動同樣優越。
固定問題有唯一係數、融合存在全域解、班表存在最小值，
都不能取代時間外驗證與營運資料校準。
後續調度仍需核實逐班剩餘容量、分鐘到達形狀、完整增班公里、
起站行駛與整備時間，及實際車隊和行車限制。

\appendix
\section{有限生成錐封閉的完整證明}\label{app:cone}

本附錄證明引理 \ref{lem:cone}。
不能只說「封閉集合的線性像封閉」，因為該敘述一般不成立；
此處需要有限生成的特性。

設 $C=\{Pw:w\ge0\}$。顯然 $0\in C$；
任意兩點的非負線性組合仍屬 $C$，所以 $C$ 是凸錐。
取 $z\in C$ 的一組非負表示 $z=Pw$，只看正權重的欄向量。
若它們線性相依，存在非零 $h$ 支持在這些欄上，使 $Ph=0$。
必要時令 $h$ 換號，保證至少一個 $h_j>0$。定義
\begin{equation}
t=\min_{j:h_j>0}\frac{w_j}{h_j}>0,\qquad w'=w-th.
\end{equation}
則 $w'\ge0$ 且 $Pw'=Pw$，至少一個原本的正係數變為 0，
並未新增原支持以外的欄。重複有限次，
就能以線性獨立的欄向量非負組合表示 $z$；$z=0$ 亦可取空組合。

因此 $C$ 是有限多個由線性獨立欄生成的錐之聯集。
考慮其中一個滿 column rank 的矩陣 $A$。
若 $z_r=A\theta_r\to z$ 且 $\theta_r\ge0$，則
\begin{equation}
\theta_r=(A^\top A)^{-1}A^\top z_r
\longrightarrow\theta=(A^\top A)^{-1}A^\top z\ge0.
\end{equation}
由連續性 $A\theta=\lim_rA\theta_r=z$，故此錐封閉。
空組合生成的 $\{0\}$ 也封閉。
有限個封閉集合的聯集封閉，所以 $C$ 封閉，證畢。

\section{OLS、Lasso 與變數分析的適用界線}

這些方法曾用於專案比較或資料探索，並非目前保留的三個核心方法全都同時使用。
本附錄補充解存在性的關係與推論條件，避免把探索指標當成必經的篩選規則。

\subsection{OLS 與 Lasso 的解}
對完整設計矩陣 $X$（含截距），OLS 解
$\min_\beta\norm{y-X\beta}^2$ 一定存在：
$\operatorname{col}(X)$ 是有限維封閉線性子空間，
$y$ 對此空間的正交投影唯一存在。
但係數唯一須有 $\rank X=p$；rank 不足時可能有多個係數對應相同擬合。

Lasso \cite{lasso} 解
\[
\min_{\beta_0,\beta}\frac1n\norm{y-\beta_0\one-X\beta}^2
+\lambda\norm\beta_1,\qquad\lambda>0.
\]
消去最適截距後，目標對 $\beta$ 連續且至少為 $\lambda\norm\beta_1$，
故次水平集有界且封閉，保證取得最小值。
L1 懲罰不是嚴格凸，相關或重複欄位可能使係數解不唯一。
這與 Ridge 用正 L2 懲罰保證唯一的機制不同；
不存在「Lasso 理論上一定比 Ridge 更適合」的結論。

\subsection{相關、VIF 與交互作用}
Spearman 相關是對排名取 Pearson 相關；有相同值時使用平均秩。
它衡量單調關聯，不證明因果，也不能由低相關排除非線性或條件交互作用。
VIF 對非恆定欄位 $X_j$ 的定義為
\[
\operatorname{VIF}_j=\frac1{1-R_j^2},
\]
其中 $R_j^2$ 來自以其他欄位與截距迴歸 $X_j$。
在固定且滿 column rank 的設計、獨立同變異誤差的 OLS 中，
\[
\Var(\widehat\beta_j\mid X)
=\frac{\sigma^2}{\sum_i(X_{ij}-\overline X_j)^2}\operatorname{VIF}_j.
\]
這是 $(X^\top X)^{-1}$ 的分割矩陣表示，
說明共線性放大係數估計變異；不是「VIF 超過某數就必須刪除」的定理。
加入車站與時段的主效果、交互作用時應處理參考組，
避免完整 dummy 與截距的精確線性相依。
Ridge 與 GAM 的正懲罰可處理 rank 不足的最小化，
但不會憑空創造能辨認因果效果的獨立資訊。

\subsection{F 檢定與 AIC 不可無條件套用}
對兩個巢狀 OLS 模型，完整模型 rank 為 $p_F$，
簡化模型比它少 $q>0$ 個自由參數，且 $n>p_F$，傳統統計量為
\[
F=\frac{(\operatorname{RSS}_R-\operatorname{RSS}_F)/q}
{\operatorname{RSS}_F/(n-p_F)}.
\]
在虛無假設成立、固定設計且誤差為 $N(0,\sigma^2I)$ 時，
投影差與完整殘差是互相正交的 Gaussian 分量，
各自平方和除以 $\sigma^2$ 為自由度 $q$ 與 $n-p_F$ 的獨立卡方，
故比例服從 $F_{q,n-p_F}$。
同一活動內各站、各時段相關時，$\sigma^2I$ 假設不自動成立；
不能把全部 504 列當成獨立重複而直接沿用此精確分布。
經資料選模後再做相同檢定，也需要處理選擇造成的推論偏差。

AIC 的形式為 $-2\log L(\widehat\theta)+2k$ \cite{aic}，
是在指定概似與正則條件下估計資訊損失的準則，
不是保證未來預測最準的定理。
一般 $k$ 不可任意替代懲罰模型的有效自由度；
本案目前的超參數比較主要依時間向前驗證，並非用 AIC 統一選所有模型。

\section{調度公式的可重算教學例子}

此例完全是假設資料，用來驗算式子，不是北門最佳班表的實驗結果。
單方向 0--10 分鐘、每分鐘到達 2 人、起始無排隊。
原班在第 10 分鐘提供 20 人容量；候選增班在第 3、5、7 分鐘，
各容量 20 人、各成本 10 元，最小班距 3 分鐘，提前量均已符合。
採 $v=196/60$。

由於每一班都能清空當時候車，對任一相鄰服務間隔長度 $\Delta$，
排隊面積為 $\int_0^\Delta2t\,dt=\Delta^2$。
故不增班為 $10^2=100$；只在第 5 分鐘增班為 $5^2+5^2=50$；
在第 3、7 分鐘增班為 $3^2+4^2+3^2=34$。

\begin{table}[htbp]
\centering\small
\caption{有限候選調度的完整可行方案列舉}
\begin{tabular}{@{}lrrrr@{}}
\toprule
增班時刻 & 班數 & $W$（人分鐘） & 成本（元） & $\Delta C+vW$（元）\\
\midrule
無 &0&100&0&326.67\\
3 &1&58&10&199.47\\
5 &1&50&10&173.33\\
7 &1&58&10&199.47\\
3、7 &2&34&20&131.07\\
\bottomrule
\end{tabular}
\end{table}
3 與 5、5 與 7 都只間隔 2 分鐘，所以不能同選。
這個有限集合的最佳值是 $20+(196/60)\times34=131.0666\ldots$。
強制兩班仍有解，強制三班則無解，
因定理 \ref{thm:dispatch} 的非空前提不成立。
上述列舉與既有批次候車函數、MILP 教學驗算一致；
數值核對只檢查這個例子的算術與表示，不取代理論證明。

\section{程式、輸出與資料來源索引}\label{app:implementation}

本文依本機專案 2026 年 10 月 18 日檢視狀態編寫，
當時 Git 基準版本為 \texttt{5a331a1}，另含尚未提交的理論與營運文件。
下列路徑相對於專案根目錄，可供核對本文公式與目前實作。

\begin{longtable}{@{}p{73mm}p{86mm}@{}}
\toprule
位置 & 本文用途\\
\midrule
\endfirsthead
\toprule 位置 & 本文用途\\\midrule
\endhead
\path{src/forecasting/linear.py} & Ridge 的標準化、$n\lambda$ 懲罰與非負預測處理。\\
\path{src/forecasting/gam_fit.py} & GAM 的候選、基底與擬合介面。\\
\path{src/forecasting/gam_stable.py} & 係數差分、全係數懲罰及數值檢查。\\
\path{src/forecasting/weighted_tree.py} & 北門散場重疊權重、場內正規化與 XGBoost 設定。\\
\path{src/queueing.py} & 有限期流體到達、批次上車與候車面積；不含本文完整增班最佳化。\\
\path{result/forecast/validation.json} & 既有實作核對與資料使用界線；不是新活動驗證。\\
\path{README.md}、\path{process.md} & 現行保留結果、模型理由與執行方式。\\
\path{docs/metro_operations.md} & 官方時刻表、班距、容量與調度時間來源；平均班距與安全門檻須區分。\\
\path{result/queue_dispatch_v1/beimen_holiday_baseline.csv} & 已整理的官方假日基準班表；不是已最佳化班表。\\
\path{docs/theory_paper/} & 本文件的可編輯 LaTeX 原始碼與證明。\\
\bottomrule
\end{longtable}

\clearpage
\begin{thebibliography}{99}
\setlength{\itemsep}{5pt}

\bibitem{ridge}
Hoerl, A. E., \& Kennard, R. W. (1970).
Ridge regression: Biased estimation for nonorthogonal problems.
\emph{Technometrics}, 12(1), 55--67.
\href{https://doi.org/10.1080/00401706.1970.10488634}{DOI: 10.1080/00401706.1970.10488634}.
用途：Ridge 正則化背景；本文第 3 節依本案損失尺度給出矩陣推導。

\bibitem{gam}
Hastie, T., \& Tibshirani, R. (1986).
Generalized additive models.
\emph{Statistical Science}, 1(3), 297--310.
\href{https://doi.org/10.1214/ss/1177013604}{DOI: 10.1214/ss/1177013604}.
用途：加法模型的基本形式；不代表本文所有基底與門檻皆取自原文。

\bibitem{pspline}
Eilers, P. H. C., \& Marx, B. D. (1996).
Flexible smoothing with B-splines and penalties.
\emph{Statistical Science}, 11(2), 89--102（正文；討論續至 121 頁）.
\href{https://doi.org/10.1214/ss/1038425655}{DOI: 10.1214/ss/1038425655}.
用途：B-spline 係數差分懲罰；本文另加入全係數 L2 穩定項。

\bibitem{xgb}
Chen, T., \& Guestrin, C. (2016).
XGBoost: A scalable tree boosting system.
\emph{Proceedings of KDD}, 785--794.
\href{https://arxiv.org/abs/1603.02754}{作者論文：arXiv:1603.02754}.
用途：第 2.2 節正則化目標、二階展開、最佳葉值及分裂增益；
本文將其寫成加權平方損失的特例。

\bibitem{breiman}
Breiman, L. (1996).
Stacked regressions.
\emph{Machine Learning}, 24, 49--64.
\href{https://doi.org/10.1007/BF00117832}{DOI: 10.1007/BF00117832}.
用途：以交叉驗證預測學習非負線性組合的經典依據；
不保證任意資料都優於最佳單模型。

\bibitem{pystacked}
Ahrens, A., Hansen, C. B., \& Schaffer, M. E. (2023).
pystacked: Stacking generalization and machine learning in Stata.
\emph{The Stata Journal}, 23(4), 909--931.
\href{https://doi.org/10.1177/1536867X231212426}{DOI: 10.1177/1536867X231212426}.
用途：第 2 節 stacking 與限制方式；包含非負但不強制總和為 1 的選項。

\bibitem{bertsekas}
Bertsekas, D. P. (2009).
\emph{Convex Optimization Theory}. Athena Scientific.
\href{https://web.mit.edu/dimitrib/www/convexduality.html}{作者教材與相關資源（MIT）}.
用途：有限維凸分析、有限生成錐的封閉性；
本文附錄 A 提供本案所需的完整獨立證明。

\bibitem{boyd}
Boyd, S., \& Vandenberghe, L. (2004).
\emph{Convex Optimization}. Cambridge University Press.
\href{https://web.stanford.edu/~boyd/cvxbook/}{作者開放教材（Stanford）}.
用途：第 4--5 章凸最佳化與 KKT；
本文明確另證可行性、存在性，未僅由凸性推論有解。

\bibitem{fpp}
Hyndman, R. J., \& Athanasopoulos, G. (2021).
\emph{Forecasting: Principles and Practice}, 3rd ed. OTexts.
\href{https://otexts.com/fpp3/tscv.html}{第 5.10 節：時間序列交叉驗證}.
用途：只使用過去資料的向前驗證；活動分群與年度切分為本案設計。

\bibitem{cawley}
Cawley, G. C., \& Talbot, N. L. C. (2010).
On over-fitting in model selection and subsequent selection bias in performance evaluation.
\emph{Journal of Machine Learning Research}, 11, 2079--2107.
\href{https://www.jmlr.org/papers/v11/cawley10a.html}{JMLR 原文}.
用途：模型選擇與反覆評估可能造成的樂觀偏差。

\bibitem{little}
Little, J. D. C. (1961).
A proof for the queuing formula: $L=\lambda W$.
\emph{Operations Research}, 9(3), 383--387.
\href{https://doi.org/10.1287/opre.9.3.383}{DOI: 10.1287/opre.9.3.383}.
用途：長期排隊量、到達率與時間的關係；
本文第 7 節另證有限視野面積恆等式，未直接套用穩態假設。

\bibitem{urban}
Larson, R. C., \& Odoni, A. R. (1981).
\emph{Urban Operations Research}. Prentice-Hall.
\href{https://web.mit.edu/urban_or_book/www/book/chapter4/4.4.html}{MIT 開放教材，第 4.4 節}.
用途：排隊系統計數與 Little 關係的面積解釋。

\Needspace{6\baselineskip}
\bibitem{pan}
Pan, H., Liu, Y., Liu, Z., \& Hu, H. (2022).
Passenger inflow control along rail corridor for recurrent heavy passenger flow mitigation.
\emph{Transportation Research Record}, 2676(9), 59--75.
\href{https://doi.org/10.1177/03611981221085527}{DOI: 10.1177/03611981221085527}.
用途：共同考慮等待與營運成本、MIP 與 MPC 的研究方向；
本次核對摘要，本文式 \eqref{eq:dispatch} 是自行設計，非原文公式重現。

\bibitem{mol}
勞動部（2026）。基本工資／最低工資之制訂與調整經過。
\href{https://www.mol.gov.tw/1607/28690/2282/2284/2290/7121/}{勞動部官方資料}.
用途：2026 年 1 月 1 日起最低時薪 196 元；
換算乘客時間價值是本案研究假設，非該頁提出的交通估價方法。

\bibitem{lasso}
Tibshirani, R. (1996).
Regression shrinkage and selection via the lasso.
\emph{Journal of the Royal Statistical Society: Series B}, 58(1), 267--288.
\href{https://doi.org/10.1111/j.2517-6161.1996.tb02080.x}{DOI: 10.1111/j.2517-6161.1996.tb02080.x}.
用途：附錄 B 的 L1 懲罰背景；存在性論證在本文直接給出。

\bibitem{aic}
Akaike, H. (1974).
A new look at the statistical model identification.
\emph{IEEE Transactions on Automatic Control}, 19(6), 716--723.
\href{https://doi.org/10.1109/TAC.1974.1100705}{DOI: 10.1109/TAC.1974.1100705}.
用途：附錄 B 的 AIC 形式與資訊準則背景。
\end{thebibliography}

\end{document}
